% !TeX root = document.tex
\section*{Вступ}
\textbf{Актуальність теми}. Розвиток сучасних інформаційних технологій та систем штучного інтелекту нерозривно пов'язаний з цифровою обробкою зображень. У таких високоточних сферах, як медична візуалізація, аналіз супутникових знімків, астрофізика та комп'ютерногий зір, точність передачі та обробки графічних даних має вирішальне значення. Операція масштабування є базовою процедурою в цих галузях, проте її виконання класичними методами завжди пов'язане з ризиком втрати, викривлення або зміщення первинної інформації.

Більшість базових методів масштабування балансує між продуктивністю та якістю, що на практиці виражається у деградації різких контурів. Спроби вирішити цю проблему шляхом впровадження класичної глобальної інтерполяції історично стикалися зі значними обчислювальними бар'єрами та структурною вразливістю до крайових спотворень, відомих як феномен Рунге. З огляду на це, традиційні алгоритми не здатні повною мірою задовольнити вимоги сучасних задач, для яких ключовими критеріями є водночас обчислювальна стійкість та еталонна точність перетворень.

Головним завданням даної роботи є створення нового методу, здатного нівелювати зазначені недоліки та зберегти високу структурну точність пікселів, що набуває особливого значення при сильному стисненні зображень. Перспективним шляхом вирішення цієї актуальної задачі є застосування інтерполяційних формул за вузлами Чебишева, які широко використовуються в обчислювальній математиці і можуть бути ефективно адаптовані для алгоритмів зміни масштабу зображень.

\textbf{Мета роботи}: розробити та програмно реалізувати новий метод масштабування цифрових зображень на основі другої барицентричної інтерполяційної формули на основі вузлів Чебишева.

Для досягнення поставленої мети необхідно вирішити наступні задачі:

\begin{enumerate}
	\item Ознайомитися з існіючими методами масштабування зображень.
	\item Розробити програмне забезпечення для масштабування зображень використовуючи новий запропонований та стандартні реалізовані методи.
	\item Провести тестування та порівняльний аналіз результатів роботи розробленого алгоритму з класичними методами комп'ютерної графіки.
\end{enumerate}

\textbf{Об'єкт дослідження} — алгоритми та процеси масштабування комп'ютерної графіки.

\textbf{Предмет дослідження} — математичний апарат барицентричної інтерполяції за вузлами Чебишева в задачах зміни розміру зображень.

\textbf{Методи дослідження}. 

\textbf{Наукова новизна отриманих результатів}.

\textbf{Практичне значення отриманих результатів}. Теоретичні положення роботи доведено до рівня конкретного прикладного рішення — створено та оптимізовано алгоритм масштабування цифрових зображень мовою C++. Розроблений математично строгий інструмент може бути застосований у спеціалізованих задачах комп'ютерного зору (наприклад, у медичній візуалізації або аналізі супутникових знімків), де вимоги до відсутності артефактів та безкомпромісної структурної точності даних (особливо при сильному зменшенні масштабу) є критично вищими за вимоги до швидкості виконання обчислень.
\newpage
